% ============================================================
%  Příprava na maturitu – Mechatronika
%  Kompiluj: pdflatex maturita_mechatronika.tex
%  Kódování: UTF-8
% ============================================================

\documentclass[a4paper,10pt]{article}

% --- Balíčky ---
\usepackage[T1]{fontenc}
\usepackage[utf8]{inputenc}
\usepackage[czech]{babel}
\usepackage[left=2.4cm, right=1.8cm, top=1.8cm, bottom=1.8cm]{geometry}
\usepackage{lmodern}
\usepackage{microtype}
\usepackage{parskip}
\usepackage{titlesec}
\usepackage{enumitem}
\usepackage{xcolor}
\usepackage{hyperref}
\usepackage{tabularx}
\usepackage{booktabs}
\usepackage{array}
\usepackage{tikz}
\usepackage{eso-pic}
\usepackage{mdframed}
\usepackage{amsmath}
\usepackage{amssymb}
\usepackage{pgfplots}
\pgfplotsset{compat=1.18}
\usetikzlibrary{arrows.meta, positioning, shapes.geometric, calc}

% --- Barvy ---
\definecolor{accent}{HTML}{03254E}
\definecolor{stripeblue}{HTML}{568EA3}
\definecolor{medgray}{HTML}{888888}
\definecolor{lightblue}{HTML}{EAF2F8}
\definecolor{lightyellow}{HTML}{FDFBE4}
\definecolor{lightgreen}{HTML}{EAF7EA}

% --- Modrý pruh vlevo ---
\AddToShipoutPictureBG{%
  \begin{tikzpicture}[remember picture, overlay]
    \fill[stripeblue]
      (current page.north west)
      rectangle
      ([xshift=10mm] current page.south west);
    \fill[accent!60]
      ([xshift=10mm] current page.north west)
      rectangle
      ([xshift=12mm] current page.south west);
  \end{tikzpicture}%
}

\hypersetup{colorlinks=true, urlcolor=accent, linkcolor=accent}

\titleformat{\section}
  {\large\bfseries\color{accent}}
  {\thesection.}
  {0.5em}
  {}
  [\color{accent}\titlerule]

\titlespacing{\section}{0pt}{16pt}{6pt}

\newmdenv[
  backgroundcolor=lightblue,
  linecolor=stripeblue,
  linewidth=1pt,
  roundcorner=4pt,
  innerleftmargin=8pt,
  innerrightmargin=8pt,
  innertopmargin=6pt,
  innerbottommargin=6pt,
  skipabove=4pt,
  skipbelow=4pt
]{pojmybox}

\newmdenv[
  backgroundcolor=lightyellow,
  linecolor=accent!40,
  linewidth=1pt,
  roundcorner=4pt,
  innerleftmargin=8pt,
  innerrightmargin=8pt,
  innertopmargin=6pt,
  innerbottommargin=6pt,
  skipabove=4pt,
  skipbelow=8pt
]{vysvetlenibox}

\newmdenv[
  backgroundcolor=lightgreen,
  linecolor=accent!30,
  linewidth=1pt,
  roundcorner=4pt,
  innerleftmargin=8pt,
  innerrightmargin=8pt,
  innertopmargin=6pt,
  innerbottommargin=6pt,
  skipabove=4pt,
  skipbelow=8pt
]{vzorcebox}

\pagestyle{plain}

% ============================================================
\begin{document}

% --- Záhlaví ---
\begin{minipage}[t]{0.75\textwidth}
    {\Huge\bfseries\color{accent} Příprava na maturitu}\\[4pt]
    {\large\color{stripeblue} Mechatronika}\\[2pt]
    \textcolor{medgray}{\small Suchánek Lukáš \quad SPŠ Prosek \quad 2026}
\end{minipage}
\hfill
\begin{minipage}[t]{0.22\textwidth}
    \raggedleft
    \begin{tikzpicture}
        \draw[color=stripeblue, rounded corners=4pt, line width=1pt]
            (0,0) rectangle (3,1.2)
            node[midway, align=center, color=accent, font=\small\bfseries]
            {Otázek: 22};
    \end{tikzpicture}
\end{minipage}

\vspace{6pt}
\noindent\color{accent}\rule{\linewidth}{1.5pt}
\vspace{2pt}

% ============================================================
\section{Diskrétní řízení -- použití, signály, diskretizace, PSD regulátor}

\begin{pojmybox}
\textbf{Klíčové pojmy:}
diskrétní signál, vzorkování, Shannonův teorém, kvantování,
A/D převodník, perioda vzorkování $T_s$, Z-transformace,
PSD regulátor, dopředná/zpětná diference, antiwindup
\end{pojmybox}

\begin{vysvetlenibox}
\textbf{Proč diskrétní řízení?}\\
Moderní řídicí systémy (PLC, MCU) jsou digitální -- zpracovávají signál
\textit{nespojitě}, v pravidelných časových okamžicích.
Spojitý signál musíme nejprve \textbf{vzorkovat} a \textbf{kvantovat}.

\textbf{Proces převodu spojitého signálu na diskrétní:}
\begin{enumerate}[noitemsep, leftmargin=1.4em, topsep=4pt]
    \item \textbf{Vzorkování} -- získání hodnoty signálu v časových okamžicích $t = k \cdot T_s$
    \item \textbf{Držení} -- hodnota se udržuje konstantní po dobu převodu (sample \& hold)
    \item \textbf{Kvantování} -- zaokrouhlení napětí na nejbližší diskrétní úroveň
    \item \textbf{Kódování} -- přiřazení binárního kódu kvantizační úrovni
\end{enumerate}

\textbf{Shannonův vzorkovací teorém}

\begin{vzorcebox}
$$f_s \geq 2 \cdot f_{max}$$
Vzorkovací frekvence musí být alespoň \textbf{dvojnásobek} nejvyšší
frekvence v signálu. Jinak dojde k \textbf{aliasingu} (překrývání spekter).
\end{vzorcebox}

\textbf{Příklad aliasingu:}\\
Pokud vzorkujeme signál 50\,Hz frekvencí 80\,Hz ($< 2 \cdot 50$),
výsledek bude vypadat jako signál s frekvencí $80 - 50 = 30$\,Hz.
Tento falešný signál se nazývá \textit{alias}.

\textbf{Antialiasingový filtr}\\
Abychom předešli aliasingu, řadíme před A/D převodník \textbf{dolní propust},
k která potlačí frekvence vyšší než $f_s / 2$. Typicky se používá
Besselův nebo Butterworthův filtr 2.--4. řádu.

\textbf{Kvantování a rozlišení A/D převodníku}

\begin{vzorcebox}
$$\Delta U = \frac{U_{ref}}{2^n - 1}$$
kde $n$ je počet bitů A/D převodníku, $\Delta U$ je kvantizační krok.\\
Příklad: 12-bitový ADC s $U_{ref} = 3{,}3\,\text{V}$:
$\Delta U = \frac{3{,}3}{4095} \approx 0{,}806\,\text{mV}$
\end{vzorcebox}

\textbf{Typy A/D převodníků:}
\begin{itemize}[noitemsep, leftmargin=1.4em, topsep=4pt]
    \item \textbf{Flash ADC} -- nejrychlejší (ns), mnoho komparátorů,
          pro vysoké frekvence (osciloskopy, rádiové přijímače)
    \item \textbf{SAR ADC} (Successive Approximation) -- střední rychlost ($\mu$s),
          nízká spotřeba; nejpoužívanější v mikrokontrolérech
    \item \textbf{Sigma-Delta ADC} -- vysoké rozlišení (24 bitů), pomalejší;
          pro audio, přesná měření (vážní snímače, teploměry)
    \item \textbf{Dual-slope ADC} -- velmi přesný, pomalý;
          pro multimetry a měřicí přístroje
\end{itemize}

\textbf{Parametry A/D převodníků:}
\begin{itemize}[noitemsep, leftmargin=1.4em, topsep=4pt]
    \item \textbf{Rozlišení} -- počet bitů ($n$); určuje počet úrovní $2^n$
          (8 bitů = 256, 10 bitů = 1024, 12 bitů = 4096, 16 bitů = 65\,536)
    \item \textbf{Vzorkovací frekvence} -- maximální počet převodů za sekotu
          (kS/s, MS/s); např. Arduino Uno: 10-bitový ADC, 15\,kS/s
    \item \textbf{Referenční napětí $U_{ref}$} -- určuje rozsah měření;
          může být interní (např. 1,1\,V, 2,5\,V) nebo externí
    \item \textbf{Kvantizační chyba} -- $\pm \frac{1}{2}$ LSB (Least Significant Bit);
          nelze odstranit, je vlastností převodu
    \item \textbf{Nelinearita} -- INL (Integral Non-Linearity) a DNL (Differential);
          udává odchylku od ideální přímky
    \item \textbf{Čas převodu} -- doba mezi startem převodu a platným výstupem
\end{itemize}

\textbf{Zapojení A/D převodníku v obvodu:}
\begin{itemize}[noitemsep, leftmargin=1.4em, topsep=4pt]
    \item \textbf{Vstupní impedance} -- ADC zatěžuje zdroj; vysoká impedance zdroje
          způsobuje chyby (doporučuje se buffer/OP)
    \item \textbf{Filtr na vstupu} -- RC dolní propust před ADC potlačí vysoké frekvence
          a sníží aliasing i šum
    \item \textbf{Stínění} -- analogové vedení vést odděleně od digitálních signálů;
          zemní plocha jako stínění
\end{itemize}

\textbf{D/A převodníky}\\
Pro převod digitálního signálu na analogový (např. řízení motoru, generování signálů):

\textbf{Princip funkce:}
\begin{itemize}[noitemsep, leftmargin=1.4em, topsep=4pt]
    \item \textbf{R-2R žebřík} -- rezistorová síť, každý bit přispívá proudem
          podle své váhy; MSB má největší vliv, LSB nejmenší
    \item \textbf{PWM + filtr} -- pulzně šířková modulace s RC filtrem;
          jednoduché, levné, pro nenáročné aplikace (stmívání LED, řízení otáček)
    \item \textbf{DAC s proudovým výstupem} -- přesné zdroje proudu řízené bity;
          vyžaduje převodník proudu na napětí (OP)
\end{itemize}

\textbf{Parametry D/A převodníků:}
\begin{itemize}[noitemsep, leftmargin=1.4em, topsep=4pt]
    \item \textbf{Rozlišení} -- počet bitů; určuje minimální změnu výstupního napětí
    \item \textbf{Usazovací čas (settling time)} -- doba, za kterou výstup dosáhne
          cílové hodnoty s danou přesností (typicky 1--10\,$\mu$s)
    \item \textbf{Glitch} -- přechodový jev při změně kódu (např. 0111 $\to$ 1000);
          potlačuje se sample\&hold obvodem na výstupu
    \item \textbf{Reference} -- interní nebo externí; určuje rozsah výstupního napětí
\end{itemize}

\textbf{Použití D/A převodníků:}
\begin{itemize}[noitemsep, leftmargin=1.4em, topsep=4pt]
    \item \textbf{Generování signálů} -- sinus, trojúhelník, pila pro audio,
          testovací generátory
    \item \textbf{Řízení pohonů} -- analogový řídicí signál pro frekvenční měniče,
          servozesilovače
    \item \textbf{Zdroje napětí} -- programovatelné laboratorní zdroje,
          kalibrátory měřicích přístrojů
\end{itemize}

\textbf{PSD regulátor -- diskrétní analogie PID}\\
Derivace se nahradí \textbf{diferencí}, integrál \textbf{sumou}:

\begin{vzorcebox}
$$u_k = K_p \cdot e_k
+ K_i \cdot T_s \sum_{j=0}^{k} e_j
+ K_d \cdot \frac{e_k - e_{k-1}}{T_s}$$
kde $e_k = w_k - y_k$ je regulační odchylka v kroku $k$,
$T_s$ je perioda vzorkování.
\end{vzorcebox}

\textbf{Implementační detaily PSD regulátoru:}
\begin{itemize}[noitemsep, leftmargin=1.4em, topsep=4pt]
    \item \textbf{Proporcionální složka} -- okamžitá reakce na chybu
    \item \textbf{Sumační člen} -- akumulace chyby; pozor na přetečení
    \item \textbf{Derivační člen} -- predikce budoucího vývoje;
          v praxi se často filtruje (šumová citlivost)
    \item \textbf{Pozicionální algoritmus} -- výpočet absolutní hodnoty $u_k$
    \item \textbf{Inkrementální algoritmus} -- výpočet změny $\Delta u_k$;
          bezpečnější při přepnutí do manuálního režimu
\end{itemize}

\textbf{Vliv periody vzorkování $T_s$}

\begin{center}
\begin{tabularx}{0.95\linewidth}{@{} l X X @{}}
    \toprule
    & \textbf{Příliš velké $T_s$} & \textbf{Příliš malé $T_s$} \\
    \midrule
    Stabilita     & systém může kmitat nebo být nestabilní & dobrá \\
    Výpočetní zátěž & nízká & vysoká \\
    Aliasing      & hrozba & nehrozí \\
    Doporučení    & $T_s < \frac{T_{63\%}}{10}$
                  & zbytečné plýtvání výkonem \\
    \bottomrule
\end{tabularx}
\end{center}

\textbf{Volba periody vzorkování v praxi:}
\begin{itemize}[noitemsep, leftmargin=1.4em, topsep=4pt]
    \item \textbf{Teplotní soustavy} -- pomalé ($T_{63\%} \approx 10$\,s): $T_s = 1$\,s
    \item \textbf{Hladiny, tlaky} -- střední ($T_{63\%} \approx 1$\,s): $T_s = 100$\,ms
    \item \textbf{Pohony, poloha} -- rychlé ($T_{63\%} \approx 100$\,ms): $T_s = 10$\,ms
    \item \textbf{Elektronické zatížení} -- velmi rychlé: $T_s = 10$\,--\,100\,$\mu$s
\end{itemize}

\textbf{Antiwindup}\\
Antiwindup je ochrana proti přetečení integrační složky při saturaci
akčního zásahu. Když regulátor dosáhne maxima (např. 100\,\% výkonu),
integrace se zastaví -- jinak by se "naintegrovala" velká hodnota
a výstup by pozdě přestál přeběh.

\begin{itemize}[noitemsep, leftmargin=1.4em, topsep=4pt]
    \item \textbf{Podmíněná integrace} -- integrace běží, pouze když $|e_k| < \text{limit}$
    \item \textbf{Back-calculation} -- rozdíl mezi ideálním a saturovaným výstupem
          se vrací na vstup integrátoru se zpožděním
    \item \textbf{Clamping} -- zastavení integrace při dosažení limitu
\end{itemize}

\textbf{Z-transformace}\\
Pro analýzu diskrétních systémů se používá Z-transformace:
$$X(z) = \sum_{k=0}^{\infty} x_k \cdot z^{-k}$$
Převod mezi spojitým a diskrétním světem:
\begin{itemize}[noitemsep, leftmargin=1.4em, topsep=4pt]
    \item \textbf{Dopředná diference} -- $s \approx \frac{z-1}{T_s}$
    \item \textbf{Zpětná diference} -- $s \approx \frac{1-z^{-1}}{T_s}$
    \item \textbf{Tustinova metoda} -- $s \approx \frac{2}{T_s} \cdot \frac{z-1}{z+1}$;
          zachovává stabilitu, nejpoužívanější
\end{itemize}
\end{vysvetlenibox}

% ============================================================

% ============================================================
\end{document}
